\newpage
\lesson{21/10}
\chapter{Quantum Field Theory}\\

Up to this point the constructed field theory was purely classical:
now the goal is to quantize it.

There are several ways to do so, but the simplest to approach is
\textbf{Canonical Quantization}, while the better and more difficult approach,
\textbf{Path Integral Quantization} will be employed for interacting theories
in the Foundations of QFT course.
\section{Canonical Quantization: general idea}
The core idea is to promote the classical fields to \textit{field operators}
$$\phi(x^\mu) \longrightarrow \hat \phi(x^\mu),$$
specify its commutation relations with other fields and then express
the field as a \textit{Fourier expansion}. The final step is to compute
the Hamiltonian of the system and find its eigenstates and eigenvalues.

Schematically, the process can be divided in
\begin{itemize}
    \item Write the Lagrangian of the system and compute the Hamiltonian of the
        system;
    \item Promote the field(s) to operator(s) and set the commutation relations
        among the fields;
    \item Expand the field(s) as Fourier Transforms;
    \item Solve for the eigenvalues and eigenstates of the Hamiltonian.
\end{itemize}
\subsection{Canonical Quantization of the Scalar Field}
The Lagrangian of the scalar field is
$$\mathcal{L}(\phi(x^\mu),\pd_\mu\phi(x^\mu)) =
\frac{1}{2}\pd_\mu\phi\pd^\mu\phi - \frac{1}{2}m^2\phi^2:$$
consider now the \textit{canonical conjugated field}, defined with a parallelism
to conjugated momenta in classic Hamiltonian mechanics
$$\pi(x^\mu):=\pdv{\mathcal{L}}{\pd_0\phi}=\pdv[]{\mathcal{L}}{\dot\phi}=
\dot\phi.$$
Again in analogy with Hamiltonian Mechanics, the Hamiltonian of the system is
$$\mathcal{H} := \pi\dot\phi-\mathcal{L} = \frac{1}{2}\dot\phi^2 - \frac{1}{2}
\pd_i\phi\pd^i\phi+\frac{1}{2}m^2\phi^2$$
\textbf{Note 1} The Hamiltonian formulation \textit{loses} the manifest Lorentz
covariance.\\
\textbf{Note 2} This Hamiltonian is definite positive, as is expected from a
non-interacting theory.

Now $\phi$ and $\pi$ are promoted to operators, and their commutator is
\textbf{defined at fixed time and different space}:
$$\left[\hat \phi(t,\mathbf{x}),\hat \pi(t,\mathbf{y})\right]=
i\hbar\delta^3(\mathbf{x-y})$$
again in analogy to the quantization of Hamiltonian mechanics for point-like
particles. There is an explanation for the result of this commutator, connected
with Noether charges that will not be covered in the lectures, but is present
in the notes of prof. Boglione.

The other commutation relations are trivial
$$\left[\hat \phi(t,\mathbf{x}),\hat\phi(t,\mathbf{y})\right]=0,\;
\left[\hat\pi(t,\mathbf{x}),\hat\pi(t,\mathbf{y})\right]=0$$
but have a deep implication: if a measurement is taken at the same time, but
in different places, their result will be \textbf{independent}, salvaging
causality in the theory.

Consider now the Fourier Transform of $\hat\phi$:
$$\hat\phi(x^\mu) = \frac{1}{(2\pi)^\frac{3}{2}}\int \dd ^4k \delta(k^2-m^2)
\vartheta(k^0)\left[\hat A(\mathbf{k})e^{-ikx}+\hat A^\dagger(\mathbf{k})e^{ikx}\right]
$$
where $k^2=k_\mu k^\mu$ and $kx = k_\mu x^\mu$\\
\textbf{Notes}
\begin{itemize}
    \item The measure of integration is Lorentz-invariant;
    \item $k^\mu$ plays the role of the 4-momentum, so $k^0=E_k$;
    \item The Heaviside $\vartheta$ guarantees $k^0=E_k\geq0$;
    \item The Dirac $\delta$ guarantees the mass-shell dispersion
        relations and fixes one of the components of $k^\mu$: for
        convention, $k^0$ is used.
\end{itemize}
Using the $\delta$, an integration over $k^0$ is performed, recalling
$$\delta(f(k))=\sum_{a}^{}\frac{\delta(k-a)}{|f'(a)|}$$
where $a$ are the simple zeroes of $f(k)$, the argument of the $\delta$
in the F.T. has poles at
$$k^2=m^2\iff k^0 =\pm\sqrt{\mathbf{k}^2+m^2}$$
and so
$$\delta(k^2-m^2) = \frac{\delta(k^0-\sqrt{\mathbf{k}^2+m^2})}{2k_0} +
\frac{\delta(k^0+\sqrt{\mathbf{k}^2+m^2})}{2|k^0|}$$
the second contribution is not admissible, since the $\vartheta$ filters out
the $k^0<0$ values:
$$\delta(k^2-m^2)\vartheta(k^0) = \frac{\delta(k^0-\omega_k)}{2\omega_k},\;
\text{with }\omega_k=+\sqrt{\mathbf{k}^2+m^2}$$
The F.T. becomes
$$\hat\phi(x^\mu) = \int \frac{\dd^3k}{(2\pi)^\frac{3}{2}}\frac{1}{2\omega_k}
\left[\hat A(\mathbf{k})e^{-ikx}+\hat A^\dagger(\mathbf{k})e^{ikx}\right]$$
and, defining
$$\hat a^{(+)} = \frac{\hat A^{(+)}}{\sqrt{2\omega_k}}$$
one gets
$$\boxed{\hat\phi(x^\mu) = \int \frac{\dd^3k}{(2\pi)^\frac{3}{2}}\frac{1}
{\sqrt{2\omega_k}}\left[\hat a(\mathbf{k})e^{-ikx}+\hat a^\dagger(\mathbf{k})e^{ikx}
\right]}$$
and the conjugated field
$$\boxed{\hat\pi(x^\mu) = \hat{\dot\phi} =
\int \frac{\dd^3k}{(2\pi)^\frac{3}{2}}\frac{1}
{\sqrt{2\omega_k}}i\omega_k\left[-\hat a(\mathbf{k})e^{-ikx}+\hat a^\dagger(\mathbf{k})
e^{ikx}\right].}$$

Since now all the information on the field is encoded in $\hat a^{(+)}$, it is
necessary to find out what are these operators: consider the commutation relation
$$\left[\hat\phi(t,\mathbf{x}),\hat\pi(t,\mathbf{y})\right] = 
\int \frac{\dd^3k}{(2\pi)^{\flatfrac{3}{2}}(2\omega_k)^{\flatfrac{1}{2}}}
\int \frac{i\omega_{k'}\dd^3k'}{(2\pi)^{\flatfrac{3}{2}}(2\omega_{k'})^
{\flatfrac{1}{2}}}\left[(\hat a_ke^{-ikx}+a_k^\dagger e^{ikx}),
(-a_{k'}e^{-ik'y}+a^\dagger_{k'}e^{ik'y})\right]$$
it is NOT easy to show that for this commutation relation to yield
$i\delta^3(\mathbf{x}-\mathbf{y})$, the $a$ operators must obey
$$\boxed{
    \left[\hat a_k,\hat a_{k'}\right]=0,\;
    \left[\hat a_k^\dagger,\hat a_{k'}^\dagger\right]=0,\;
    \left[\hat a_k,\hat a_{k'}^\dagger\right]=\delta^3(\mathbf{k-k'}).
}$$
these operators obey the \textbf{same algebra of the creation and annihilation
operators for the quantum harmonic oscillator} and they are interpreted
as operators that \textbf{create and destroy local ripples in the fields}, which
are \textbf{the particles we see.}

Now, the Hamiltonian is computed with the F.Ts of the fields
\begin{gather*}
H =\int \dd ^3x \mathcal{H} =  \frac{1}{2}\int \dd ^3x \left(
\hat\pi^2 +\pd_i\hat\phi\pd^i\hat\phi + m^2\hat\phi^2\right)=\\
=\frac{1}{2}\int \dd^3x \int \frac{\dd^3k\dd^3k'}{(2\pi)^3\sqrt{2\omega_k}
\sqrt{2\omega_{k'}}}\left[
-(\omega_k\omega_{k'}+\mathbf{k\cdot k'})(\hat a_ke^{-ikx}-\hat a_k^\dagger e^{ikx})
(\hat a_{k'}e^{-ik'x}-\hat a^\dagger_{k'}e^{ik'x})+\right. \\
\left. +m^2(\hat a_ke^{-ikx}+\hat a_k^\dagger 
e^{ikx})(\hat a_{k'}e^{-ik'x}+\hat a_{k'}^\dagger e^{ikx})\right].
\end{gather*}
Integrating over $x$ only touches the 3-vector part of the exponentials,
yielding Dirac $\delta^3(\mathbf{k\pm k'})$ and greatly simplifying matters.
$$H =[...]= \frac{1}{2}\int \dd^3k \omega_k(\hat a_k\hat a^\dagger_k+
\hat a^\dagger_k\hat a_k)$$
The commutation relation with $\mathbf{k=k'}$ yields a problematic result
$$\left[\hat a(k),\hat a^\dagger(k)\right] = \delta^3(0)\to\infty$$

Using the commutation relation in H yields
$$H = \int \dd^3k \omega_k\left(\hat a_k \hat a_k^\dagger+\frac{1}{2}
\delta^3(0)\right):$$
the form of this Hamiltonian density is in the form of that of the familiar
quantum harmonic oscillator, with the key difference that the $0$-point energy
is $\delta^3(0)$; this enforces the idea that the scalar field can
be though of as an infinite sum of Harmonic Oscillators,
so its ground state energy is the divergent series
$$\sum_{i=0}^{\infty}\frac{1}{2}\omega.$$

This problem with the $\delta$ can be resolved by setting the \textbf{0-point
energy to 0}, since this is always allowed in physics; to do so, the Hamiltonian
must be redefined so that the annihilation operators are on the right: this is
called \textbf{normal ordering}
$$N(H) = :H: =\int \dd^3k \omega_k \hat a^\dagger_k\hat a_k$$
and it consists in putting the annihilation operators, $\hat a_k$ in this case,
on the right side of the Hamiltonian density $\mathcal{H}$.
